\documentclass[aspectratio=169, 10pt]{beamer}

% =============================================================================
% CẤU HÌNH BIÊN DỊCH & GÓI LỆNH (TƯƠNG THÍCH HOÀN HẢO OVERLEAF / TEXLIVE)
% =============================================================================
\usepackage{iftex}
\ifPDFTeX
  \usepackage[utf8]{vietnam}
\else
  \usepackage{fontspec}
  \setsansfont{DejaVu Sans}
  \setmonofont{DejaVu Sans Mono}
\fi

\usepackage{amsmath,amssymb}
\usepackage{booktabs}
\usepackage{colortbl}
\usepackage{tabularx}
\usepackage{graphicx}
\usepackage{xcolor}
\usepackage{tikz}
\usetikzlibrary{positioning,arrows.meta}
\usepackage{tcolorbox}
\tcbuselibrary{skins}
\usepackage{hyperref}

% Đường dẫn tìm ảnh linh hoạt (hỗ trợ cả chạy từ thư mục gốc, thư mục slides/, hoặc drive)
\graphicspath{{./}{slides/}{../}{../slides/}{drive-download-20260925T212557Z-1-001/}{../drive-download-20260925T212557Z-1-001/}}

% =============================================================================
% ĐỊNH NGHĨA HỆ MÀU SẮC CHUYÊN NGHIỆP (VINSOC THEME)
% =============================================================================
\definecolor{socblue}{RGB}{18, 48, 92}
\definecolor{socteal}{RGB}{14, 116, 144}
\definecolor{socsky}{RGB}{2, 132, 199}
\definecolor{socdark}{RGB}{30, 41, 59}
\definecolor{socbg}{RGB}{248, 250, 252}
\definecolor{socborder}{RGB}{226, 232, 240}
\definecolor{socgreen}{RGB}{22, 101, 52}
\definecolor{socorange}{RGB}{194, 65, 12}
\definecolor{socred}{RGB}{185, 28, 28}

% =============================================================================
% CẤU HÌNH GIAO DIỆN BEAMER & LỀ TRANG
% =============================================================================
\usetheme{default}
\usefonttheme{professionalfonts}
\setbeamertemplate{navigation symbols}{}
\setbeamersize{text margin left=6mm, text margin right=6mm}

% Màu sắc các phần tử chính của Beamer
\setbeamercolor{structure}{fg=socblue}
\setbeamercolor{palette primary}{bg=socblue,fg=white}
\setbeamercolor{palette secondary}{bg=socteal,fg=white}
\setbeamercolor{palette tertiary}{bg=socdark,fg=white}
\setbeamercolor{item}{fg=socblue}
\setbeamercolor{subitem}{fg=socteal}
\setbeamercolor{block title}{bg=socblue, fg=white}
\setbeamercolor{block body}{bg=socbg, fg=socdark}
\setbeamercolor{block title example}{bg=socteal, fg=white}
\setbeamercolor{block body example}{bg=socbg, fg=socdark}

% Tiêu đề Slide (Frametitle) - Thanh ngang 1 dòng thanh thoát, tinh gọn
\makeatletter
\setbeamerfont{frametitle}{size=\normalsize, series=\bfseries}
\setbeamercolor{frametitle}{bg=socblue, fg=white}
\setbeamertemplate{frametitle}{%
  \nointerlineskip%
  \begin{beamercolorbox}[wd=\paperwidth,ht=3.2ex,dp=1.2ex,left]{frametitle}%
    \hspace*{2.5ex}\usebeamerfont{frametitle}\insertframetitle%
    \ifx\insertframesubtitle\@empty%
    \else%
      \hskip 2em plus 1fill{\footnotesize\color{socsky!45!white}\textmd{\insertframesubtitle}}\hspace*{2.5ex}%
    \fi%
  \end{beamercolorbox}%
  \vspace{-1pt}%
  {\color{socsky}\hrule height 1.2pt}%
}
\makeatother

% Chân trang Slide (Footline) - Rõ ràng, dễ đọc, cân đối hài hòa
\setbeamercolor{footline}{bg=socdark, fg=white}
\setbeamertemplate{footline}{%
  {\color{socsky!60!black}\hrule height 0.8pt}%
  \leavevmode%
  \hbox{%
    \begin{beamercolorbox}[wd=.33\paperwidth,ht=3.2ex,dp=1.2ex,left]{footline}%
      \hspace*{2.5ex}\fontsize{6}{7}\selectfont\textbf{\color{socsky}VinSOC} \textcolor{white!35}{|} \textmd{\color{white!85}Security Operations Center}
    \end{beamercolorbox}%
    \begin{beamercolorbox}[wd=.44\paperwidth,ht=3.2ex,dp=1.2ex,center]{footline}%
      \footnotesize\color{white!90}Tuần 4: Bộ sinh bất thường \& benchmark
    \end{beamercolorbox}%
    \begin{beamercolorbox}[wd=.23\paperwidth,ht=3.2ex,dp=1.2ex,right]{footline}%
      \footnotesize\color{white!80}Trang \;\textbf{\color{white}\insertframenumber} / \inserttotalframenumber\hspace*{2.5ex}
    \end{beamercolorbox}%
  }%
  \vskip0pt%
}

% =============================================================================
% MACRO & THÀNH PHẦN HỘP TINH GỌN (CHỐNG TRÀN VERTICAL TUYỆT ĐỐI)
% =============================================================================
% Hộp Nhận xét nghiệp vụ ngắn gọn
\newcommand{\nhanxet}[1]{%
  \vspace{0.12em}%
  \begin{tcolorbox}[
    enhanced,
    colback=socsky!6,
    colframe=socsky,
    boxrule=0pt,
    leftrule=3pt,
    arc=1mm,
    left=6pt, right=6pt, top=3pt, bottom=3pt,
    fontupper=\scriptsize
  ]
    \textbf{\color{socsky!90!black}Nhận xét nghiệp vụ:} #1
  \end{tcolorbox}%
}

% Hộp Chú thích biểu đồ tinh gọn
\newcommand{\chuthichbieudo}[1]{%
  \vspace{0.12em}%
  \begin{tcolorbox}[
    enhanced,
    colback=socbg,
    colframe=socborder,
    boxrule=0.6pt,
    leftrule=3pt,
    borderline west={3pt}{0pt}{socteal},
    arc=1mm,
    left=6pt, right=6pt, top=3pt, bottom=3pt,
    fontupper=\scriptsize
  ]
    \textbf{\color{socteal}Ý nghĩa:} #1
  \end{tcolorbox}%
}

% Hộp thẻ card nội dung gọn nhẹ (thay thế beamer block to cồng kềnh)
\newtcolorbox{soccard}[2][]{%
  enhanced,
  colback=socbg,
  colframe=socborder,
  boxrule=0.6pt,
  leftrule=2.5pt,
  borderline west={2.5pt}{0pt}{socblue},
  arc=1mm,
  left=5pt, right=5pt, top=3.5pt, bottom=3.5pt,
  fontupper=\scriptsize,
  title={\textbf{\color{socblue}#2}},
  coltitle=socblue,
  attach title to upper={\par\vspace{1.5pt}},
  #1
}

% Lệnh chèn ảnh an toàn hỗ trợ đa đường dẫn
\newcommand{\chenanh}[2][width=\linewidth]{%
  \IfFileExists{#2}{%
    \includegraphics[#1]{#2}%
  }{%
    \IfFileExists{slides/#2}{%
      \includegraphics[#1]{slides/#2}%
    }{%
      \IfFileExists{drive-download-20260925T212557Z-1-001/#2}{%
        \includegraphics[#1]{drive-download-20260925T212557Z-1-001/#2}%
      }{%
        \IfFileExists{../drive-download-20260925T212557Z-1-001/#2}{%
          \includegraphics[#1]{../drive-download-20260925T212557Z-1-001/#2}%
        }{%
          \centering\fbox{\textit{\color{gray}[Tệp hình ảnh: #2]}}\par%
        }%
      }%
    }%
  }%
}

% Cấu hình bảng biểu tinh gọn
\setlength{\tabcolsep}{4pt}
\renewcommand{\arraystretch}{1.18}

% =============================================================================
% THÔNG TIN BÁO CÁO
% =============================================================================
\title{BỘ SINH BẤT THƯỜNG TỔNG HỢP VÀ BENCHMARK}
\subtitle{Tuần 4: Tiêm nhãn, đánh giá và tinh chỉnh UEBA}
\author{\textbf{VinSOC Security Operations Center}}
\institute{Trung Tâm Vận Hành An Ninh Mạng VinSOC}
\date{Tháng 10/2026}

% =============================================================================
% BẮT ĐẦU TÀI LIỆU SLIDE
% =============================================================================
% Biên dịch: XeLaTeX (khuyến nghị) hoặc pdfLaTeX, chạy hai lần.
% Nội dung nguồn: reports/week4/tom_tat_tuan4.md (bản chạy lại commit ff3cf9c)
% Số liệu: experiments/grid/dev/ (ngày 36--42, 5 run dev_seed20261036..40)
%          experiments/grid/test/ (ngày 43--60, 5 run test_seed20261061..65)
% Hình: reports/week4/hinh/ (chép từ experiments/grid/{dev,test}/figures/)
\newcommand{\duongdan}[1]{{\scriptsize\texttt{\detokenize{#1}}}}
\setbeamerfont{block title}{size=\normalsize,series=\bfseries}
\setbeamerfont{block body}{size=\small}
\hypersetup{pdftitle={UEBA Tuần 4 - Bộ sinh bất thường và benchmark},pdfauthor={VinSOC Security Operations Center}}
\begin{document}

\begin{frame}[plain]
\begin{tikzpicture}[remember picture,overlay]
\fill[socblue] (current page.north west) rectangle ([yshift=-24pt]current page.north east);
\fill[socsky] ([yshift=-24pt]current page.north west) rectangle ([yshift=-27pt]current page.north east);
\end{tikzpicture}
\begin{center}
{\LARGE\bfseries\color{socblue}KHUNG BENCHMARK UEBA}\\[.6em]
{\large\bfseries\color{socteal}Tuần 4: Bộ sinh bất thường tổng hợp và benchmark}\\[.7em]
{\normalsize Tiêm nhãn vào dữ liệu đánh giá, so sánh mô hình với baseline}\\[1em]
{\small Isolation Forest \quad LOF \quad One-Class SVM}\\[1.5em]
{\small VinSOC Security Operations Center}\\[.4em]
{\footnotesize Tháng 10/2026}
\end{center}
\end{frame}

\begin{frame}{01. Nội dung công việc tuần 4}
\begin{columns}[T,onlytextwidth]
\begin{column}{.48\textwidth}
\begin{block}{Bộ sinh và quy trình tiêm}
\begin{itemize}\setlength\itemsep{7pt}
\item 6 kịch bản: 4 theo đề cương, 2 bổ sung.
\item Tiêm đúng 1\% số dòng User theo tài khoản/ngày.
\item Tiêm ở tầng log, rồi tính lại 41 đặc trưng.
\item Nhãn và manifest cho từng run.
\end{itemize}
\end{block}
\end{column}
\begin{column}{.48\textwidth}
\begin{exampleblock}{Benchmark}
\begin{itemize}\setlength\itemsep{7pt}
\item Chia train, dev, test theo thời gian (mục 5.1).
\item 3 mô hình, 3--6 cấu hình mỗi mô hình, lặp 5 seed.
\item So sánh với 4 mốc tham chiếu.
\item Chọn cấu hình trên dev, chạy test một lần.
\end{itemize}
\end{exampleblock}
\end{column}
\end{columns}
\vspace{.8em}
\nhanxet{Trên test (ngày 43--60), One-Class SVM đạt PR-AUC \textbf{0,182 $\pm$ 0,023}, gấp \textbf{7 lần} luật đếm số lần thất bại và \textbf{18 lần} ngẫu nhiên.}
\end{frame}

\begin{frame}{02. Chia train, dev và test theo thời gian}
\normalsize
\begin{tabularx}{\linewidth}{@{}p{2cm}p{2.8cm}X@{}}
\toprule\rowcolor{socblue!10}\textbf{Khối} & \textbf{Ngày} & \textbf{Vai trò}\\\midrule
\textbf{Train} & 1--42 (70\%) & Dựng hồ sơ hành vi, kho khuôn sự kiện và huấn luyện mô hình cho test.\\[10pt]
\textbf{Dev} & 36--42 & Nằm \textbf{trong train}: mô hình, hồ sơ và luật loại trừ dựng từ ngày 1--35. Dùng để chọn cấu hình.\\[10pt]
\textbf{Test} & 43--60 (30\%) & Đánh giá một lần với cấu hình đã chốt, không tinh chỉnh bằng nhãn test.\\\bottomrule
\end{tabularx}
\vspace{1em}
\begin{itemize}\setlength\itemsep{5pt}
\item Train giữ nguyên trong mọi run tiêm, chỉ ngày đánh giá bị tiêm.
\item Imputer và scaler chỉ \texttt{fit} trên train.
\item Một run tiêm gồm \textbf{một khối dev hoặc test và một seed}.
\end{itemize}
\end{frame}

\begin{frame}{03. Nguyên tắc dựng sự kiện tổng hợp}
\normalsize
\begin{enumerate}\setlength\itemsep{10pt}
\item \textbf{Dùng sự kiện thật làm khuôn.} Ưu tiên sự kiện train của nạn nhân, thiếu thì dùng tài khoản cùng loại.
\item \textbf{Thay đổi theo hành vi kịch bản.} Điều chỉnh thời gian, máy nguồn/đích và các trường cần thiết như kết quả xác thực hoặc LogonType.
\item \textbf{Ghi vào bản sao log đánh giá.} Log gốc và train giữ nguyên.
\item \textbf{Tính lại đặc trưng sau tiêm.} Trích lại toàn bộ 41 đặc trưng, gồm lịch sử 7 ngày.
\item \textbf{Nhãn lưu riêng.} Log đưa vào trích đặc trưng không có cột đánh dấu tiêm.
\end{enumerate}
\end{frame}

\begin{frame}{04. Hồ sơ hành vi và điều kiện chọn nạn nhân}
\begin{columns}[T,onlytextwidth]
\begin{column}{.48\textwidth}
\begin{block}{Hồ sơ học từ train của khối}
\begin{itemize}\setlength\itemsep{7pt}
\item Máy nguồn và máy đích quen thuộc.
\item LogonType thường dùng.
\item Giờ hoạt động điển hình.
\item Khoảng nghỉ dài nhất của tài khoản.
\item Ngưỡng khoá mô phỏng $L=5$.
\end{itemize}
\end{block}
\end{column}
\begin{column}{.48\textwidth}
\begin{exampleblock}{Điều kiện chung}
\begin{itemize}\setlength\itemsep{7pt}
\item Chỉ tài khoản người dùng (User).
\item Có ít nhất 7 ngày hoạt động trong train.
\item Mỗi tài khoản bị tiêm tối đa một lần trong một run.
\item Loại dòng tài khoản/ngày mà luật ngưỡng đã cảnh báo trên log gốc.
\end{itemize}
\end{exampleblock}
\end{column}
\end{columns}
\vspace{.6em}
\small Mỗi kịch bản còn có điều kiện riêng để hành vi tiêm khác với lịch sử của nạn nhân.
\end{frame}

\begin{frame}{05. Sáu kịch bản bất thường tổng hợp}
\small
\begin{tabularx}{\linewidth}{@{}p{4.4cm}X@{}}
\toprule\rowcolor{socblue!10}\textbf{Kịch bản} & \textbf{Mô tả hành vi được tiêm}\\\midrule
\textbf{Brute-force} & 8--25 lần đăng nhập thất bại từ một máy nguồn lạ trong 5--30 phút, sau $L=5$ lần sai thì bị khoá.\\[7pt]
\textbf{Password spraying} & Một máy nguồn lạ thử sai 1--3 lần trên 5--13 tài khoản, \textbf{cả chiến dịch trong một cửa sổ 30--90 phút}.\\[7pt]
\textbf{Đăng nhập ngoài giờ} & \textbf{Dời toàn bộ một ngày hoạt động điển hình} của người làm ban ngày sang 0--7 giờ sáng.\\[7pt]
\textbf{Bùng nổ máy trạm mới} & Đăng nhập thành công từ 4--10 máy nguồn chưa từng dùng trong train, trong 0,5--2 giờ.\\[7pt]
\textbf{Ngủ đông thức dậy} & Tài khoản im lặng lâu hơn khoảng nghỉ dài nhất của chính nó bỗng có một ngày hoạt động điển hình (giữ giờ gốc).\\[7pt]
\textbf{Đổi kiểu đăng nhập} & \textbf{Cả một ngày} đăng nhập chuyển sang LogonType chưa từng dùng; ưu tiên chiều \textbf{2 $\rightarrow$ 5/3} theo đề cương, phần còn lại 3 $\rightarrow$ 10/2.\\\bottomrule
\end{tabularx}
\end{frame}

\begin{frame}{06. Bốn chỗ đã sửa sau khi đối chiếu đề cương mục 5.2}
\small
\begin{tabularx}{\linewidth}{@{}p{2.6cm}>{\raggedright\arraybackslash}X>{\raggedright\arraybackslash}X@{}}
\toprule\rowcolor{socblue!10}\textbf{Kịch bản} & \textbf{Bản trước} & \textbf{Bản hiện tại}\\\midrule
\textbf{Spraying} & Mỗi nạn nhân một cửa sổ riêng, chiến dịch trải 9--23 giờ (lỗi code) & Một cửa sổ chung 30--90 phút, chiến dịch $\geq$ 5 nạn nhân\\[6pt]
\textbf{Ngoài giờ} & Thêm 10--40 sự kiện đêm vào ngày đã có hoạt động ban ngày, tín hiệu bị pha loãng & Dời cả ngày sang đêm\\[6pt]
\textbf{Đổi LogonType} & Thêm 3--10 sự kiện vào ngày có hàng trăm sự kiện (tín hiệu khoảng 2\%), ngược chiều đề cương & Cả ngày đổi loại logon, ưu tiên 2 $\rightarrow$ 5/3\\[6pt]
\textbf{Ngủ đông} & Ngày phát lại bị dời giờ ngẫu nhiên (lỗi code) & Giữ nguyên giờ trong ngày\\\bottomrule
\end{tabularx}
\vspace{.6em}
\nhanxet{ROC-AUC của OCSVM trên test: ngoài giờ tăng \textbf{0,49 $\rightarrow$ 0,91}, đổi LogonType tăng \textbf{0,44 $\rightarrow$ 0,99}. Hai kịch bản này trước đây khó phát hiện là do cách tiêm, không phải do mô hình.}
\end{frame}

\begin{frame}{07. Quy trình tiêm và đánh giá}
\centering
\begin{tikzpicture}[
  task/.style={draw=socblue!55, fill=socbg, rounded corners=2pt,
    line width=.8pt, align=center, text=socdark,
    font=\small, text width=3.6cm, minimum height=1.1cm, inner sep=5pt},
  flow/.style={-{Latex[length=2.6mm,width=1.7mm]}, line width=1.05pt, draw=socteal}
]
\node[task] (profile) at (0,2.1) {\textbf{1. Hồ sơ hành vi}\\Học từ train};
\node[task] (rule) at (4.5,2.1) {\textbf{2. Loại cảnh báo sẵn}\\Luật ngưỡng chấm log gốc};
\node[task] (template) at (9,2.1) {\textbf{3. Kho khuôn}\\Gom sự kiện thật từ train};
\node[task,fill=socteal!10,draw=socteal,text width=5.4cm] (inject) at (4.5,.2)
  {\textbf{4. Chọn và sinh}\\Chọn nạn nhân, sinh sự kiện theo 6 kịch bản};
\node[task] (write) at (0,-1.8) {\textbf{5. Ghi run}\\Ghi log đã tiêm, nhãn và manifest};
\node[task] (features) at (4.5,-1.8) {\textbf{6. Đặc trưng}\\Tính lại 41 đặc trưng};
\node[task] (evaluate) at (9,-1.8) {\textbf{7. Đánh giá}\\Benchmark trên khối đánh giá};
\draw[flow] (profile.south) -- (0,1.02) -| ([xshift=-1.8cm]inject.north);
\draw[flow] (rule.south) -- (inject.north);
\draw[flow] (template.south) -- (9,1.02) -| ([xshift=1.8cm]inject.north);
\draw[flow] (inject.south) -- ++(0,-.35) -| (write.north);
\draw[flow] (write.east) -- (features.west);
\draw[flow] (features.east) -- (evaluate.west);
\end{tikzpicture}
\par\vspace{.6em}
\small\textbf{Một run = một khối dev hoặc test + một seed cố định.}
\end{frame}

\begin{frame}{08. Ví dụ brute-force: Quy trình tiêm}
% User1, ngày 40, n=17 là giả định minh họa, không trích từ kết quả chạy.
\fontsize{9}{11}\selectfont
\par\vspace{.4em}
\begin{tabularx}{\linewidth}{@{}>{\raggedright\arraybackslash}p{3.2cm}>{\raggedright\arraybackslash}X@{}}
\toprule\rowcolor{socblue!10}\textbf{Bước} & \textbf{Luồng thực hiện cho lần tiêm này}\\\midrule
\textbf{1. Hồ sơ hành vi} & Lấy dữ liệu của User1 trong train (ngày 1--35 với khối dev) để xác định LogHost, Source quen thuộc và LogonType thường dùng, làm cơ sở sinh log bất thường.\\[7pt]
\textbf{2. Loại cảnh báo sẵn} & Kiểm tra log gốc của User1 ở ngày 40 bằng luật cảnh báo. Nếu chưa bị cảnh báo thì giữ làm ứng viên tiêm; nếu đã bị cảnh báo thì bỏ qua.\\[7pt]
\textbf{3. Kho khuôn} & Lấy log đăng nhập thất bại có mã 4625 từ train làm mẫu. Ưu tiên log của User1; nếu thiếu thì lấy log của tài khoản cùng loại. Không có mẫu phù hợp thì bỏ qua.\\[7pt]
\textbf{4. Chọn và sinh} & Chọn User1 chưa bị tiêm trong lần chạy này. Chọn ngẫu nhiên số dòng log cần thêm trong khoảng từ 8 đến 25. Các dòng log được sinh có \textbf{LogHost lạ}, \textbf{Source quen lấy từ hồ sơ hành vi} và \textbf{LogonType 3}.\\[7pt]
\textbf{5. Ghi run} & Thêm các dòng log đã sinh vào bản sao log dev, rồi gán nhãn 1 cho User1/ngày 40. Lưu riêng log đã tiêm, file nhãn và manifest chứa seed cùng các tham số tiêm.\\\bottomrule
\end{tabularx}
\par\vspace{.4em}
\end{frame}

\begin{frame}{09. Ví dụ brute-force: Quy trình đánh giá}
\fontsize{9}{11}\selectfont
\par\vspace{.4em}
\begin{tabularx}{\linewidth}{@{}>{\raggedright\arraybackslash}p{3.2cm}>{\raggedright\arraybackslash}X@{}}
\toprule\rowcolor{socblue!10}\textbf{Bước} & \textbf{Luồng thực hiện cho lần tiêm này}\\\midrule
\textbf{6. Đặc trưng} & Lấy log dev đã tiêm để tính lại các đặc trưng hành vi của User1/ngày 40, rồi xử lý dữ liệu thiếu và chuẩn hóa bằng imputer, scaler đã học trên train.\\[24pt]
\textbf{7. Đánh giá} & Dùng mô hình đã học trên train để chấm điểm log dev đã tiêm, rồi ghép điểm với nhãn theo tài khoản/ngày. Kiểm tra thứ hạng của User1/ngày 40 và tính PR-AUC, Precision@k, Recall theo ngân sách trên toàn dev.\\\bottomrule
\end{tabularx}
\par\vspace{.4em}
\end{frame}

\begin{frame}{10. Tiêm bao nhiêu bất thường vào dữ liệu?}
\small
\begin{tabularx}{\linewidth}{@{}>{\raggedright\arraybackslash}p{4.2cm}>{\raggedright\arraybackslash}X>{\raggedright\arraybackslash}X@{}}
\toprule\rowcolor{socblue!10}\textbf{Nội dung} & \textbf{Dev (ngày 36--42)} & \textbf{Test (ngày 43--60)}\\\midrule
\textbf{Mẫu số} (dòng User gốc) & 54.656 & 152.740\\[6pt]
\textbf{Số lần tiêm} $=\lceil 1\%\cdot D/0{,}99\rceil$ & \textbf{553} (92--93 mỗi kịch bản) & \textbf{1.543} (257--258 mỗi kịch bản)\\[6pt]
\textbf{Tỷ lệ đạt} & 1,01\% ở cả 5 run & 1,01\% ở cả 5 run\\[6pt]
\textbf{Sự kiện log được thêm} & 62--130 nghìn & 157--198 nghìn\\\bottomrule
\end{tabularx}
\par\vspace{1em}
Ví dụ: thêm 25 sự kiện cho User1 trong một ngày $\rightarrow$ một dòng nhãn 1.
\end{frame}

% =============================================================================
% PHẦN BENCHMARK
% =============================================================================

\begin{frame}{11. Thiết lập benchmark}
\small
\begin{tabularx}{\linewidth}{@{}p{3.4cm}X@{}}
\toprule\rowcolor{socblue!10}\textbf{Thành phần} & \textbf{Thiết lập}\\\midrule
Lưới cấu hình & Isolation Forest 3 cấu hình (\texttt{n\_estimators}, \texttt{max\_samples}, \texttt{max\_features}); LOF 4 (scaler $\times$ \texttt{n\_neighbors}); One-Class SVM 6 (scaler, \texttt{gamma}, \texttt{n\_components}).\\[5pt]
4 mốc tham chiếu & Ngẫu nhiên; \textbf{luật số lần thất bại $\geq$ 5} (ngưỡng cố định, mục 5.3); luật 6 điều kiện; z-score toàn cục.\\[5pt]
Lặp 5 seed & Mỗi seed = một bộ dữ liệu tiêm + một \texttt{random\_state} mô hình. Báo cáo trung bình $\pm$ độ lệch chuẩn.\\[5pt]
Ngưỡng cảnh báo & Phân vị $(1-c)$ của điểm \textbf{train}, $c = 5\%$. Không dùng nhãn để đặt ngưỡng.\\[5pt]
Chỉ số (mục 5.4) & PR-AUC, ROC-AUC, P@10/50/100 (toàn khối và theo ngày), Recall@5\%, Recall top-100/ngày, tỷ lệ cảnh báo, thời gian.\\[5pt]
Chọn cấu hình & PR-AUC trung bình 5 seed cao nhất trên \textbf{dev}; test sinh tự động từ cấu hình đã chốt.\\[5pt]
Nhật ký & Mọi lần fit ghi vào \texttt{experiments/logs/experiment\_log.csv} (mục 5.6).\\\bottomrule
\end{tabularx}
\end{frame}

\begin{frame}{12. Kết quả dev: toàn bộ cấu hình}
\fontsize{6.6}{7.6}\selectfont
\renewcommand{\arraystretch}{1.08}
\begin{tabularx}{\linewidth}{@{}l>{\raggedright\arraybackslash}Xccccc@{}}
\toprule\rowcolor{socblue!10}\textbf{Mô hình} & \textbf{Cấu hình} & \textbf{PR-AUC} & \textbf{ROC-AUC} & \textbf{P@100/ngày} & \textbf{Recall@5\%} & \textbf{Tỷ lệ c.báo}\\\midrule
\rowcolor{socteal!10}\textbf{OCSVM} & \textbf{standard, $\gamma$=scale, 1000 thành phần} $\star$ & \textbf{0,169 $\pm$ 0,021} & 0,908 & 0,27 & \textbf{0,58} & 9,3\%\\
OCSVM & standard, $\gamma$=scale, 300 thành phần & 0,167 $\pm$ 0,012 & 0,906 & 0,26 & 0,57 & 9,5\%\\
OCSVM & standard, $\gamma$=0,1, 300 thành phần & 0,164 $\pm$ 0,014 & \textbf{0,916} & 0,25 & 0,56 & 8,2\%\\
OCSVM & standard, $\gamma$=0,005, 300 thành phần & 0,164 $\pm$ 0,004 & 0,904 & \textbf{0,28} & 0,54 & 10,0\%\\
OCSVM & quantile, $\gamma$=scale, 300 thành phần & 0,054 $\pm$ 0,008 & 0,848 & 0,09 & 0,31 & 7,2\%\\
OCSVM & robust, $\gamma$=scale, 300 thành phần & 0,036 $\pm$ 0,012 & 0,512 & 0,07 & 0,19 & 6,4\%\\\midrule
\rowcolor{socteal!10}\textbf{IF} & \textbf{300 cây, max\_samples=1024, max\_features=1,0} $\star$ & \textbf{0,063 $\pm$ 0,008} & 0,866 & 0,12 & 0,31 & 8,2\%\\
IF & 150 cây, max\_samples=auto, max\_features=0,5 & 0,046 $\pm$ 0,003 & 0,848 & 0,09 & 0,23 & 8,4\%\\
IF & 150 cây, max\_samples=auto, max\_features=1,0 (mặc định) & 0,043 $\pm$ 0,003 & 0,846 & 0,08 & 0,23 & 7,9\%\\\midrule
\rowcolor{socteal!10}\textbf{LOF} & \textbf{robust, k=20} $\star$ & \textbf{0,049 $\pm$ 0,005} & 0,691 & 0,13 & 0,33 & 9,1\%\\
LOF & quantile, k=20 & 0,043 $\pm$ 0,002 & 0,756 & 0,08 & 0,25 & 7,1\%\\
LOF & standard, k=100 & 0,042 $\pm$ 0,010 & 0,644 & 0,10 & 0,22 & 9,6\%\\
LOF & standard, k=20 & 0,027 $\pm$ 0,008 & 0,597 & 0,07 & 0,18 & 10,4\%\\\midrule
\textit{Z-score} & max $|z|$ bền vững trên 41 đặc trưng & 0,036 $\pm$ 0,009 & 0,727 & 0,07 & 0,13 & 6,2\%\\
\textit{Luật thất bại} & số lần thất bại/ngày $\geq$ 5 (ngưỡng cố định) & 0,027 $\pm$ 0,001 & 0,631 & 0,08 & 0,21 & 4,4\%\\
\textit{Luật 6 điều kiện} & tỷ lệ luật thoả, ngưỡng phân vị train & 0,018 $\pm$ 0,001 & 0,605 & 0,06 & 0,13 & 18,8\%\\
\textit{Ngẫu nhiên} & băm tất định theo dòng & 0,011 $\pm$ 0,001 & 0,501 & 0,01 & 0,05 & 5,0\%\\\bottomrule
\end{tabularx}
\par\vspace{.2em}

\end{frame}

\begin{frame}{13. Kết quả dev: PR-AUC theo cấu hình}
\centering
\chenanh[width=\linewidth,height=6.5cm,keepaspectratio]{hinh/dev-grid-pr-auc.png}

\end{frame}

\begin{frame}{14. Kết quả test: bảng chính}{Ngày 43--60, chạy một lần, khoảng 153.570 dòng, 1.543 nhãn dương mỗi run}
\footnotesize
\begin{tabularx}{\linewidth}{@{}>{\raggedright\arraybackslash}Xcccccc@{}}
\toprule\rowcolor{socblue!10}\textbf{Mô hình (cấu hình chốt)} & \textbf{PR-AUC test} & \textbf{PR-AUC dev} & \textbf{ROC-AUC} & \textbf{P@100/ngày} & \textbf{Recall@5\%} & \textbf{Fit (s)}\\\midrule
\rowcolor{socteal!10}\textbf{OCSVM} (standard, 1000 tp) & \textbf{0,182 $\pm$ 0,023} & 0,169 & \textbf{0,917} & \textbf{0,29} & \textbf{0,64} & 43,0\\
\textbf{IF} (n=300, ms=1024) & 0,065 $\pm$ 0,003 & 0,063 & 0,888 & 0,13 & 0,32 & 13,5\\
\textbf{LOF} (k=20, robust) & 0,047 $\pm$ 0,004 & 0,049 & 0,688 & 0,11 & 0,34 & 12,8\\\midrule
\textit{Z-score toàn cục} & 0,032 $\pm$ 0,003 & 0,036 & 0,743 & 0,06 & 0,15 & 2,1\\
\textit{Luật thất bại $\geq$ 5} & 0,025 $\pm$ 0,000 & 0,027 & 0,623 & 0,08 & 0,20 & 1,1\\
\textit{Luật 6 điều kiện} & 0,015 $\pm$ 0,001 & 0,018 & 0,536 & 0,06 & 0,13 & 1,1\\
\textit{Ngẫu nhiên} & 0,010 $\pm$ 0,000 & 0,011 & 0,503 & 0,01 & 0,05 & 1,4\\\bottomrule
\end{tabularx}
\vspace{.5em}
\begin{itemize}\setlength\itemsep{3pt}\small
\item Thứ hạng giữ nguyên giữa dev và test: \textbf{OCSVM $>$ IF $>$ LOF $>$ z-score $>$ luật thất bại $>$ luật 6 điều kiện $>$ ngẫu nhiên}.
\item Với 100 cảnh báo mỗi ngày, OCSVM có \textbf{29\%} cảnh báo đúng và bắt \textbf{33\%} nạn nhân.
\end{itemize}
\end{frame}

\begin{frame}{15. Kết quả test: Precision@k (k = 10, 50, 100)}{Tỷ lệ cảnh báo đúng trong top-$k$, trung bình 5 seed}
\footnotesize
\centering
\begin{tabular*}{\linewidth}{@{\extracolsep{\fill}}lccc@{}}
\toprule
\rowcolor{socblue!10}\textbf{Mô hình} & \textbf{P@10} & \textbf{P@50} & \textbf{P@100}\\\midrule
\rowcolor{socteal!10}\textbf{OCSVM} & 0,18 & 0,24 & \textbf{0,24}\\
\textbf{IF} & \textbf{0,30} & 0,16 & 0,15\\
\textbf{LOF} & 0,18 & 0,07 & 0,08\\\midrule
\textit{Z-score toàn cục} & 0,14 & \textbf{0,26} & 0,23\\
\textit{Luật thất bại $\geq$ 5} & 0,00 & 0,00 & 0,01\\
\textit{Luật 6 điều kiện} & 0,02 & 0,03 & 0,04\\
\textit{Ngẫu nhiên} & 0,02 & 0,01 & 0,02\\\bottomrule
\end{tabular*}
\vspace{.5em}
\begin{itemize}\setlength\itemsep{3pt}\scriptsize
\item IF có P@10 0,30, z-score có P@50 0,26, nhưng đều tụt khi $k$ tăng.
\item Z-score mạnh ở top-10/ngày (0,21) rồi giảm nhanh; \textbf{luật thất bại} gần 0 ở top đầu vì điểm cao nhất là tài khoản thật vốn hay thất bại.
\end{itemize}
\end{frame}



\begin{frame}{16. Kết quả test: PR-AUC theo cấu hình chốt}{Trung bình $\pm$ độ lệch chuẩn, 5 seed}
\centering
\chenanh[width=\linewidth,height=5.2cm,keepaspectratio]{hinh/test-grid-pr-auc.png}
\par\vspace{.3em}
\chuthichbieudo{OCSVM vượt mọi mốc tham chiếu trên cả 5 seed; IF và LOF chỉ nhỉnh hơn z-score. Đường nét đứt là mức PR-AUC của bộ chấm ngẫu nhiên.}
\end{frame}

\begin{frame}{17. Kết quả test: đường Precision--Recall}{Trung bình 5 seed, dải = $\pm$1 độ lệch chuẩn}
\begin{columns}[c,onlytextwidth]
\begin{column}{.58\textwidth}
\chenanh[width=\linewidth,height=6.3cm,keepaspectratio]{hinh/test-pr-curves.png}
\end{column}
\begin{column}{.40\textwidth}
\begin{itemize}\setlength\itemsep{6pt}\scriptsize
\item \textbf{OCSVM} giữ precision khoảng 30\% tới recall 0,35: vùng vận hành hữu ích nhất.
\item \textbf{Z-score} có đỉnh precision cao (P@10/ngày 0,21) nhưng rơi nhanh: chỉ bắt nhóm nhỏ rất lộ, chủ yếu brute-force.
\item \textbf{Rule base} chỉ đạt precision khoảng 0,02 ở đầu danh sách, cao nhất 0,09 quanh recall 0,1: điểm cao nhất là tài khoản thật vốn hay thất bại, không phải nạn nhân.
\end{itemize}
\end{column}
\end{columns}
\end{frame}

% Slide 1: Bảng dữ liệu ROC-AUC theo kịch bản
\begin{frame}{18. Kết quả test theo từng kịch bản (Bảng số liệu)}{ROC-AUC, trung bình 5 seed}
\vfill
\centering
\begin{minipage}{0.85\textwidth}
\normalsize
\renewcommand{\arraystretch}{1.3}
\begin{tabular*}{\linewidth}{@{\extracolsep{\fill}}lcccc@{}}
\toprule
\rowcolor{socblue!10}\textbf{Kịch bản} & \textbf{OCSVM} & \textbf{IF} & \textbf{LOF} & \textbf{Luật TB}\\
\midrule
Brute-force & \textbf{0,994} & 0,886 & 0,453 & 0,988\\
Đổi LogonType & \textbf{0,985} & 0,957 & 0,896 & 0,359\\
Máy trạm mới & \textbf{0,975} & 0,900 & 0,919 & 0,495\\
Ngủ đông & \textbf{0,947} & 0,911 & 0,604 & 0,513\\
Ngoài giờ & \textbf{0,906} & 0,877 & 0,730 & 0,510\\
\rowcolor{socorange!10}Spraying & 0,696 & 0,797 & 0,529 & \textbf{0,875}\\
\bottomrule
\end{tabular*}
\par\vspace{.6em}
\end{minipage}
\vfill
\end{frame}

% Slide 2: Biểu đồ trực quan hóa ROC-AUC
\begin{frame}{19. Kết quả test theo từng kịch bản (Trực quan hóa)}{ROC-AUC, trung bình 5 seed}
\vfill
\centering
\chenanh[width=0.92\linewidth,height=5.8cm,keepaspectratio]{hinh/test-scenario-roc-auc.png}
\vfill
\end{frame}
\begin{frame}{20. Độ nhạy contamination }{Cấu hình chốt, trung bình $\pm$ độ lệch chuẩn 5 seed}
\centering
\chenanh[width=\linewidth,height=6.6cm,keepaspectratio]{hinh/test-contamination-sensitivity.png}
\end{frame}

\begin{frame}{21. Độ nhạy contamination: }
\small
\begin{columns}[T,onlytextwidth]
\begin{column}{.5\textwidth}
\scriptsize
\begin{tabularx}{\linewidth}{@{}c>{\centering\arraybackslash}X>{\centering\arraybackslash}Xc@{}}
\toprule\rowcolor{socblue!10}\textbf{$c$} & \textbf{OCSVM P / R} & \textbf{IF P / R} & \textbf{Cảnh báo}\\\midrule
0,2\% & \textbf{0,29} / 0,13 & 0,13 / 0,05 & 0,5\%\\
0,5\% & 0,27 / 0,30 & 0,10 / 0,10 & 1,1\%\\
1\% & 0,22 / 0,42 & 0,09 / 0,17 & 1,9\%\\
2\% & 0,17 / 0,54 & 0,08 / 0,24 & 3,3\%\\
5\% & 0,10 / \textbf{0,71} & 0,06 / 0,36 & 7,0\%\\\bottomrule
\end{tabularx}
\end{column}
\begin{column}{.47\textwidth}
\begin{itemize}\setlength\itemsep{3pt}\scriptsize
\item Contamination \textbf{chỉ dịch ngưỡng}: PR-AUC OCSVM chỉ dao động 0,169--0,182.
\item Chọn $c$ là chọn giữa precision và recall. Ở 0,5\%: khoảng 95 cảnh báo/ngày, cứ 4 cảnh báo có 1 tấn công, bắt 30\% nạn nhân.
\item Tỷ lệ cảnh báo thực tế cao hơn $c$ khoảng 1,4--2 lần do dữ liệu trôi so với train $\rightarrow$ cần theo dõi theo ngày.
\end{itemize}
\end{column}
\end{columns}
\end{frame}

\begin{frame}{22. Vì sao One-Class SVM đưa tấn công lên đầu danh sách tốt hơn IF?}{}
\small
\textbf{ROC-AUC chênh nhẹ, nhưng đầu danh sách cảnh báo khác hẳn:}
\par\vspace{.3em}
\begin{center}\footnotesize
\begin{tabular}{@{}lccc@{}}
\toprule\rowcolor{socblue!10}\textbf{Mô hình} & \textbf{ROC-AUC} & \textbf{PR-AUC} & \textbf{Precision top 1\% (1.536 dòng)}\\\midrule
Isolation Forest & 0,886 & 0,064 & 0,10\\
\rowcolor{socteal!10}One-Class SVM & 0,915 & 0,183 & \textbf{0,31}\\\bottomrule
\end{tabular}
\end{center}
\vspace{.2em}
\begin{columns}[T,onlytextwidth]
\begin{column}{.49\textwidth}
\begin{soccard}{Tấn công thường vượt phạm vi train}
Số đặc trưng vượt 99,9\% giá trị train:
\begin{itemize}\setlength\itemsep{1pt}
\item $\geq$ 1 đặc trưng: \textbf{36\%} dòng tấn công, chỉ \textbf{2\%} dòng bình thường.
\item $\geq$ 3 đặc trưng: 41\% là tấn công (152/372).
\item Nguồn: máy mới, LogonType mới, khoá từ nguồn lạ.
\end{itemize}
\end{soccard}
\end{column}
\begin{column}{.49\textwidth}
\begin{soccard}{Cùng đo ``khác thường'', khác cách chấm}
\begin{itemize}\setlength\itemsep{1pt}
\item \textbf{IF bão hoà}: ngưỡng cắt nằm trong min--max train $\Rightarrow$ ``vượt xa'' và ``ở rìa train'' có điểm như nhau.
\item \textbf{IF pha loãng}: mỗi lần cắt chọn ngẫu nhiên 1/39 đặc trưng, trong khi tín hiệu chỉ nằm ở vài đặc trưng.
\item \textbf{OCSVM}: điểm tăng theo khoảng cách tới vùng train $\Rightarrow$ càng xa, càng bất thường.
\end{itemize}
\end{soccard}
\end{column}
\end{columns}
\end{frame}

\begin{frame}{23. Bằng chứng: IF tiêu tốn đầu danh sách vào đuôi dài của dữ liệu thật}
\footnotesize
\begin{tabularx}{\linewidth}{@{}>{\raggedright\arraybackslash}Xcc|cc@{}}
\toprule\rowcolor{socblue!10}\textbf{Số đặc trưng vượt 99,9\% train} & \textbf{Số dòng} & \textbf{Dòng tấn công} & \textbf{IF: chỗ top 1\% (đúng)} & \textbf{OCSVM: chỗ top 1\% (đúng)}\\\midrule
0 (trong phạm vi train) & 149.845 & 981 & \textcolor{socred}{\textbf{928}} (51) & 356 (165)\\
1 & 2.782 & 278 & 419 (60) & 771 (95)\\
2 & 568 & 132 & 71 (19) & 134 (69)\\
$\geq$ 3 & 372 & 152 & 118 (26) & \textbf{275 (142)}\\\bottomrule
\end{tabularx}
\par\vspace{.4em}
\begin{columns}[T,onlytextwidth]
\begin{column}{.5\textwidth}
\scriptsize
\begin{tabularx}{\linewidth}{@{}>{\raggedright\arraybackslash}Xcc@{}}
\toprule\rowcolor{socblue!10}\textbf{Tỷ lệ nạn nhân lọt top 1\%} & \textbf{IF} & \textbf{OCSVM}\\\midrule
Brute-force & 10\% & \textbf{75\%}\\
Đổi LogonType & 17\% & \textbf{46\%}\\
Máy trạm mới & 5\% & \textbf{37\%}\\
Ngủ đông & 5\% & \textbf{8\%}\\
Ngoài giờ & \textbf{4\%} & 0\%\\
Spraying & \textbf{20\%} & 16\%\\\bottomrule
\end{tabularx}
\end{column}
\begin{column}{.47\textwidth}
\begin{itemize}\setlength\itemsep{3pt}\scriptsize
\item \textbf{IF}: 60\% top 1\% là dòng trong phạm vi train, chỉ 5\% đúng. Đây chủ yếu là tài khoản thật hay bị khoá hoặc đăng nhập nhiều (\texttt{failure\_locked\_out\_share}, \texttt{log\_total\_logons}).
\item \textbf{OCSVM}: 77\% top 1\% là dòng vượt train, 26\% đúng.
\item IF chỉ hơn ở spraying và ngoài giờ, hai kịch bản gần như không vượt train $\Rightarrow$ lợi thế của OCSVM gắn với tấn công tạo giá trị cực trị.
\end{itemize}
\end{column}
\end{columns}
\end{frame}

\begin{frame}{24. Đối chiếu đề cương và giới hạn}
\begin{columns}[T,onlytextwidth]
\begin{column}{.48\textwidth}
\begin{block}{Đầu ra tuần 4}
\begin{itemize}\setlength\itemsep{4pt}\small
\item[\color{socgreen}\checkmark] Bộ sinh 4 kịch bản đề cương + 2 bổ sung.
\item[\color{socgreen}\checkmark] Tiêm 1,01\% ở cả 10 run.
\item[\color{socgreen}\checkmark] Bảng 3 mô hình $\times$ $\geq$ 3 cấu hình.
\item[\color{socgreen}\checkmark] Biểu đồ PR, độ nhạy contamination.
\item[\color{socgreen}\checkmark] 5 seed kèm độ lệch chuẩn.
\item[\color{socgreen}\checkmark] Đủ chỉ số mục 5.4, nhật ký chung.
\end{itemize}
\end{block}
\end{column}
\begin{column}{.48\textwidth}
\begin{exampleblock}{Điểm lệch và giới hạn}
\begin{itemize}\setlength\itemsep{3pt}\scriptsize
\item LOF (HNSW) và OCSVM (Nystroem + SGD) là bản xấp xỉ; OCSVM RBF chính xác không chạy được trên 315.000 dòng.
\item Z-score lấy max $|z|$ trên 41 đặc trưng.
\item Chiều 2 $\rightarrow$ 5/3 chỉ có 5--8 nạn nhân/run do log thiếu tài khoản dùng type 2.
\item Nhãn tổng hợp chỉ cho cận trên lạc quan; train có 601 dòng giống tấn công.
\item Chỉ đánh giá User; chưa có một lệnh chạy toàn bộ chuỗi.
\end{itemize}
\end{exampleblock}
\end{column}
\end{columns}
\end{frame}



\end{document}
